\section{总结与延伸}

本讲的主线可以压缩成九个连续结论：

\begin{enumerate}
  \item 早期 AI 领域由手工特征、专用术语和不同 topology 分割；
  \item 规模化深度学习先统一了端到端优化语言；
  \item seq2seq 的固定 context bottleneck 推动可微 alignment；
  \item Transformer 将 attention 提升为主要通信阶段，并与位置、残差、归一化、MLP、heads 和 masks 组成完整 package；
  \item Q/K/V 可以理解为寻址需求、可匹配标签与传输内容；
  \item nanoGPT 展示 token stream、shifted targets、causal mask、cross-entropy 与 autoregressive generation 的完整闭环；
  \item GPT、BERT 与 T5 主要由 connectivity、state source 和 objective 区分；
  \item 跨模态复用依赖 tokenization、type/position embedding 与领域 bias，而不是简单把所有问题变成文字；
  \item 表达力、可优化性、GPU 效率与规模共同产生 in-context behavior，使 prompt 能在运行时重新配置模型。
\end{enumerate}

\begin{warningbox}{最后的证据边界}
课堂证明了怎样实现一个最小 Transformer，也展示了多个领域的一手例子；它没有证明 attention 是唯一最优架构，没有证明 scaling 可无限持续，没有证明 in-context learning 等于梯度下降，也没有证明语言模型是可靠的任意程序执行器。高质量学习应把“可运行机制”“经验现象”“讲者判断”和“未来猜想”分别标记。
\end{warningbox}

\subsection{拓展阅读}

建议按问题链阅读，而不是按论文名堆叠：

\begin{enumerate}
  \item Bengio et al., \href{https://www.jmlr.org/papers/v3/bengio03a.html}{\emph{A Neural Probabilistic Language Model}}：固定窗口神经语言模型与 learned embeddings；
  \item Sutskever et al., \href{https://proceedings.neurips.cc/paper/2014/hash/a14ac55a4f27472c5d894ec1c3c743d2-Abstract.html}{\emph{Sequence to Sequence Learning with Neural Networks}}：变长 encoder--decoder；
  \item Bahdanau et al., \href{https://arxiv.org/abs/1409.0473}{\emph{Neural Machine Translation by Jointly Learning to Align and Translate}}：soft alignment 与动态 context；
  \item Vaswani et al., \href{https://arxiv.org/abs/1706.03762}{\emph{Attention Is All You Need}}：完整 Transformer package；
  \item Karpathy, \href{https://github.com/karpathy/nanoGPT}{\emph{nanoGPT}}：用最小代码把数据、模型、训练与生成接起来；
  \item Dosovitskiy et al., \href{https://arxiv.org/abs/2010.11929}{\emph{An Image is Worth 16x16 Words}}，以及 Chen et al., \href{https://arxiv.org/abs/2106.01345}{\emph{Decision Transformer}}：比较不同领域怎样定义 token；
  \item Brown et al., \href{https://arxiv.org/abs/2005.14165}{\emph{Language Models are Few-Shot Learners}} 与 Wei et al., \href{https://arxiv.org/abs/2201.11903}{\emph{Chain-of-Thought Prompting}}：区分行为现象、prompt protocol 与机制解释。
\end{enumerate}

最有效的实践顺序是：先在 Tiny Shakespeare 上逐行验证 tensor shapes 与 causal mask；再切换 subword tokenizer 和更长 context；随后分别实现 unmasked encoder 与 cross-attention；最后选择一个非文本领域，明确写出 element definition、position/type encoding、connectivity 与 objective。这样才能把“Transformer everywhere”从口号变成可验证的设计方法。

\end{document}
